%% Written 2026-11-08 from the project records.
%% Author: Kwang-Il Seon
\documentclass[english,a4paper]{my_memo}
\synctex=1
\usepackage{amsmath}
\usepackage{amssymb}
\usepackage{booktabs}
\usepackage{array}
\usepackage{url}
\urlstyle{tt}
\usepackage{babel}
\newcommand{\ev}[2]{\textnormal{\small[#1:#2]}}
\sloppy
\begin{document}

\title{Stationary finish and marching in EXHALE: why earlier molecular
LHS~1140\,b iterates were not certified}
\author{Kwang-Il Seon}
\maketitle

\section{Marching acceptance and a stationary root}

The two procedures have different stopping criteria. \textbf{Marching}
advances the time-dependent equations. In the campaign configuration,
hydrodynamics and advection use three-stage Runge--Kutta updates with
reconstructed HLLC fluxes; chemistry, carrier transport and thermal sources
also have separate implicit operations. The complete calculation is therefore
not simply a third-order Runge--Kutta integration of every term.
\ev{Main}{2308--2431,2608--2609,2777--2781}

PLM and WENO3 are spatial reconstructions, not separate stationary solvers.
Earlier initialization could switch from PLM to WENO3. Neither reconstruction
guarantees a particular final mass-flux spread; the campaign instead judges
the recorded profile variation. Every physical cell must satisfy a
maximum-to-minimum span of at most 0.004\,dex in temperature and density over
the final $K=3$ crossing times. Velocity must have a span within 1\% of its
minimum $|v|$ in cells that retain one sign and Mach number at least 0.01 in
every selected snapshot. Other cells use 1\% of $V_{\rm hs}$, the largest
$|v|$ among those cells and all selected snapshots. This is finite-band
profile acceptance, not proof that every chemical row vanishes.
\ev{M}{426--440,548--576}\ev{U}{60--61}

The checker calculates, from the final snapshot,
\begin{equation}
\tau_{\rm cross}=\max\left(
\sum_{j=j_M}^{N}\frac{\Delta r_j}{|v_j|},\;
\sum_{j=1}^{N}\frac{\Delta r_j}{c_{s,j}+|v_j|}\right),
\end{equation}
where $j_M$ is the first cell after the last cell below Mach 0.01;
$\Delta r_j$ are the widths reconstructed by the checker. If the outermost
cell is below that threshold, the first sum is infinite and coverage fails.
The selected snapshots must bracket the final window; default sampling also
requires at least 30 snapshots and no gap larger than $\tau_{\rm cross}/10$.
\ev{M}{389--418,460--465,540--564}
The preflight listed about $2.6$--$4.0\times10^5$\,s for the original five
selected seeds, not a fixed crossing time for their later evolving winds.
\ev{P}{121--128}

\textbf{The stationary finish} attempts to solve a discrete steady residual
with Jacobian-free Newton--Krylov (JFNK) iterations. Its unknowns depend on the
configuration: species can be eliminated through chemical closure, held in a
hydrodynamic solve, or included as transported unknowns. A successful
hydrodynamic Newton solve does not by itself certify every active species
balance; the final certification is a separate test.
\ev{N}{2338--2367,4271--4309,4491--4512}\ev{K/run.log}{373,872}
Marching can supply a starting state for an optional Newton finish, but
starting at a universal mass-flux spread of $10^{-3}$ is not required.

\section{What the earlier records establish}

\textbf{An atomic example.} The retained \texttt{wasp\_full\_newton}
regression is a WASP-121\,b atomic configuration. Its input requests
PLM+WENO3 and Newton. The log switches PLM to WENO3 at step 1716 and starts
the Newton finish at step 16500, with face mass-flux spread
$8.091\times10^{-2}$, not $10^{-3}$. Multiple Newton solves follow; the final
report certifies every active equation. Thus this example demonstrates a
successful atomic finish after substantial marching, not a universal
``few Newton passes'' result or a measured speed advantage for every atomic
planet. The solve carries three hydrodynamic rows; its success is not a
controlled comparison isolating the effect of molecular chemistry.
\ev{W/input.inp}{1,21,33--34}\ev{W/run.log}{1803,16661--16668,16851--16852,17001}
This v1.00 atomic regression is not one of the invalid v1.00 LHS~1140\,b
physical targets.\ev{U}{67}

\textbf{A molecular closure example.} In
\texttt{HeH1.2/xuv\_x1p28/k07}, the input sets
\texttt{Restart intent: stationary evaluate}, \texttt{Solver: Newton} and
\texttt{Coupled carrier solve: False}. The log explicitly says that no CFL
step is taken and no physical elapsed time is claimed.
\ev{K/input.inp}{32,40,42}\ev{K/run.log}{63,2970}
The worst gated carrier row, HeH$^+$ at cell 364, remains 0.996 against
$10^{-5}$ in all eight outer alternation passes. Hydrodynamic Newton returns
\texttt{info=0} in seven passes; pass 6 fails and restores the earlier
converged wind with its composition. The alternation is refused at pass 8
because the joint distance has not improved for three passes, and the final
state is not certified. This is failure of the joint stationary procedure,
not failure of every hydrodynamic Newton solve.
\ev{K/run.log}{872,1196,1517,1829,2131,2353,2442,2742,2956--2970,3080}

Such a run can finish quickly compared with a several-crossing-time march
because it takes no physical time steps and terminates without reaching a
joint root. It is not evidence that the wind physically relaxed quickly.
Historical closure outputs were used to prepare campaign seeds, but a seed
is not an accepted marching solution; in particular, the He/H~0.1 case uses
an additional init-relaxed seed.\ev{C}{26--36}\ev{Q15}{19--29}

\section{Possible difficulties with molecular chemistry and transport}\label{sec:difficulties}

These are plausible mechanisms or observations from particular diagnostics,
not five established causes of the quoted k07 outcome.

\begin{enumerate}
\item \textbf{Separated time scales.} At cell 95 in the same-state ion
record, inverse loss coefficients $n/L$ are about 2.3\,ms for H$_2^+$,
3.0\,ms for HeH$^+$ and 2.5\,s for H$_3^+$, far shorter than the seed wind
crossing times above. These are local estimates, not universal lifetimes.
Such separation can make a stationary linearization difficult, depending on
scaling and preconditioning; no Jacobian condition number or isolated causal
test is supplied by this comparison.\ev{H}{2362--2366}

\item \textbf{Trace rows and their own scales.} Carrier residuals are judged
relative to their physical terms, with a separate arithmetic-roundoff check;
small abundance alone does not establish balance.\ev{D}{1953--1969,2024--2042}
In one constructed stationary interval from the He/H~0.1 pass-9 state,
cell-500 HeH$^+$ had scale $4.35\times10^{-14}$\,cm$^{-3}$\,s$^{-1}$ and
relative residual 0.0236. Four helium rows failed, but the then-existing
constrained-cell exemption hid them. This was not a demonstration that only
one trace row refused every later interval. The corresponding marching
probe passed its rows; the exemption has since been removed.
\ev{Q}{47--66,76--82}\ev{D}{2024--2027}

\item \textbf{Admissibility boundaries.} Positivity and element budgets
limit, project or refuse candidate states. Changes of active constraints can
make the numerical residual or its sampled action nonsmooth. In the Q13b
stationary pass-9 probe, the raw Newton candidate left the helium simplex,
the projection restored feasibility, and the line search later failed.
It was not an inward direction incorrectly projected outward; the init and
marching probes behaved differently. This supports a possible boundary
obstacle, not a universal explanation.\ev{Q}{39--51}

\item \textbf{Seed-dependent chemical closure.} Earlier diagnostics in a
different molecular configuration recorded distinct chemical/front solutions
from different seeds; adding transported H$_2$ removed the small-seed
ambiguity in that test. This motivates checking root selection, but does not
establish physical bistability or two certified LHS~1140\,b wind solutions.
The low-He/H conditional alternatives also failed carrier certification.
\ev{B}{3963--3988}\ev{L}{35--44}

\item \textbf{Spatial coupling and approximate Jacobian actions.} Carrier
transport couples neighboring stationary balance rows, while local chemical
closure remains part of residual evaluation; it does not turn all chemistry
into a single global chemical solve.\ev{N}{4271--4309,4491--4512}
An exact Jacobian is linear wherever the derivative exists. A finite-difference
action can instead depend on interval, direction and discrete choices.
Earlier tests measured substantial additivity error; a later validated
private action reduced the maximum to 0.00263 over 340 products. These
records identify an approximation/evaluation issue, not intrinsic
non-additivity of the exact Jacobian or proof that transport caused the
present refusal.\ev{J}{21--31}\ev{JV}{2--12}
\end{enumerate}

\section{What marching does instead}

Marching retains the time terms and uses finite, controlled steps. Fast ions
may relax strongly within an implicit interval, but neither all chemistry
nor every carrier is forced into instantaneous equilibrium in each step.
The cell-95 record, for example, has a 0.183\,s step with both fast and slow
chemical modes.\ev{H}{2363--2366}
It is therefore a different route to a measured profile band, not a guarantee
of reaching a certified stationary root. The original preflight estimated
4.9--8.5 active host-days to the first possible three-crossing-time check for
the selected seeds; this is an estimate, not a promised completion time.
\ev{P}{121--128}
The user's decision makes marching convergence the default in v1.00 and
v1.10. A stationary root is welcome but not required; no new stationary
method or diagnostic is started merely because one is missing. The campaign
can also end NOT ACCEPTED at its finite budget.\ev{U}{60--64}

\section{The marching time step}

\textbf{The CFL number.} For an explicit discretization of a hyperbolic
system, the Courant--Friedrichs--Lewy condition requires the numerical domain
of dependence to include the physical one; it is a necessary condition, not
a sufficient stability guarantee. On a uniform Cartesian mesh, the local
Courant number is
\begin{equation}
C_j=\frac{\Delta t\,(|v_j|+c_{s,j})}{\Delta r},
\end{equation}
the fraction of a cell crossed in one step by the fastest local acoustic
signal. The input \texttt{CFL:} scales the spherical time-step prescription
below; it does not prescribe the same $C_j$ in every cell. The code comments
state $\mathrm{CFL}\le1$ using the forward-Euler SSP coefficient of one for
SSP-RK3, and assume a margin below one because the signal speeds are evaluated
from cell averages rather than reconstructed face states.
\ev{E}{20--32} This is not a general upper bound enforced by the
\texttt{CFL:} parser, which reads the supplied value; the initialized default
is 0.6. \ev{IR}{1512--1514}\ev{Par}{1022} At a fixed state, reducing this
parameter reduces the CFL candidate interval in direct proportion. The cost
and accuracy of the complete split update need not scale in that proportion,
because its state, implicit solves, and retries also matter.

The campaign marches with one global time step for all cells
(\texttt{Run mode: phys}, \texttt{CFL: 0.15}). E1a's actual k07 marching
input specifies 0.15, and the campaign's input generator retains this choice
as its default. \ev{E1}{42--44}\ev{MI}{59--60,96--102}\ev{CD}{451}
There was a recorded reason for the choice: the E1a design cites an earlier
hot-Uranus experiment in which a numerical oscillation was reported at
CFL 0.60 and 0.30 but suppressed at 0.15. That historical result motivated
using 0.15; it was not a CFL-convergence or optimization study of k07 or of
the present campaign states. The design leaves a higher-CFL comparison to a
separate experiment. \ev{ED}{243--250} The hydrodynamic CFL candidate
for each evolved cell $j=1,\ldots,N$ is
\begin{equation}
\Delta t_j^{\rm CFL}=\mathrm{CFL}\,\frac{V_j}{\max\left(A_{j-1/2}S_{j-1/2},\,
A_{j+1/2}S_{j+1/2}\right)},\qquad
S_f=\max\left(|v|+c_s\right)\ \text{over the two cells of face } f,
\end{equation}
with $A=r^2$ and $V_j=(r_{j+1/2}^3-r_{j-1/2}^3)/3$. On a locally uniform
Cartesian mesh with equal face speeds this reduces to
$\mathrm{CFL}\,\Delta r/(|v|+c_s)$. The global CFL candidate is
$\min_j\Delta t_j^{\rm CFL}$; retries can reduce the interval actually
accepted. The code uses cell-average sound speeds, with the state-dependent
adiabatic index when the caloric mixture is active, rather than the full
reconstructed Riemann signal speeds. This is a CFL prescription for the
hydrodynamic update, not a stability guarantee for the complete split
chemistry and energy update. Both the cell geometry and the signal speed
matter; the smallest cell width alone does not determine the step.
\ev{E}{12--32,44--71}\ev{Main}{2080--2118,3309--3326}

By default, every 20 accepted physical steps the code compares one full step
with two half steps from the same entry state. A resolved error above the
allowed tolerance rejects the whole interval, restores the entry state, and
retries with a smaller interval. The comparison accounts for the inner
nonlinear errors; an unresolved estimate is reported and does not by itself
reject the step. Thus sampling this estimate does not ensure that every
accepted step has a resolved temporal-error bound. The adopted trajectory is
the full-step trajectory, not the two-half-step result.
\ev{A}{1497--1515,1704--1708}\ev{Main}{2127--2139,3377--3478}

The following means are measured from the latest chunks committed by
13:34 KST on November 8, 2026. Each is the difference between consecutive
committed physical clocks divided by 50\,000 accepted steps. The journal's
\texttt{steps=50001} is the terminal loop index; the chunk logs report
50\,000 accepted steps, with rejected attempts excluded from that count.
\ev{Main}{3483--3488}
\begin{center}
\begin{tabular}{rrrrl}
\toprule
He/H & Chunk & Mean $\Delta t$ [s] & Steps / $\tau_{\rm cross}$ & Clock records \\
\midrule
14   & 49 & 0.13457 & $2.69\times10^6$ & \ev{J14}{289,295} \\
1.2  & 41 & 0.16066 & $1.60\times10^6$ & \ev{J12}{241,247} \\
2.09 & 38 & 0.16636 & $2.41\times10^6$ & \ev{J209}{223,229} \\
0.1  & 17 & 0.15609 & $1.67\times10^6$ & \ev{J01}{97,103} \\
\bottomrule
\end{tabular}
\end{center}
The last numerical column is a projection using the seed crossing times
$3.622$, $2.568$, $4.007$, and $2.602\times10^5$\,s, respectively,
not a measurement of the crossing time of the current profiles.
\ev{P}{121,125--128} With those seed times and these mean intervals, three
crossing times would take $4.8$--$8.1\times10^6$ steps; a chunk advances
1.86--3.13\% of a crossing time. The intervals and crossing times can change
as the wind evolves.

\section{Local (cell-by-cell) time stepping}

EXHALE offers \texttt{Time stepping: Local}, which uses each cell's own
CFL interval instead of the global minimum. Among the recognized run modes,
it is accepted only with \texttt{Run mode: init}; with
\texttt{Run mode: phys} the run stops at startup. Unequal intervals multiply
the flux on the two sides of an internal face by different amounts, so their
contributions to the total material inventory need not cancel. The update is
a relaxation iterate, not an advance along one common physical clock.
\ev{E}{72--81}\ev{IR}{1083--1090,2715--2730}

For the archived He/H~1.2 snapshot at chunk 41,
$t=334184.60265$\,s, a reconstruction using a constant adiabatic index
$5/3$, the recorded pressure and mass density, and midpoint faces gives a
largest-to-smallest CFL-interval ratio of 707.8 and a median-to-smallest
ratio of 14.10 over the 500 physical cells. The smallest interval occurs at
cell 50, $r=1.00957\,R_p$, and the largest at cell 500. The density column
is converted to mass density with the code's fixed hydrogen-atom mass,
$1.67353284\times10^{-24}$\,g. These ratios are an estimate under the
constant-$\gamma$ assumption, not the intervals calculated with the active
caloric mixture, and they do not bound the actual speed-up of the complete
relaxation. \ev{S12}{2--4,33--536}\ev{IO}{52}\ev{G}{209--210}\ev{E}{44--49}

\textbf{Possible advantages.}
\begin{itemize}
\item The large outer cells can advance through larger relaxation intervals
than the cells near the base, which may accelerate some stationary
relaxations. The speed-up for these campaign states has not been measured.
\item A locally stepped relaxation may provide a useful starting state for
a global-step march or an optional Newton finish; reaching a suitable state
is not guaranteed.
\end{itemize}

\textbf{Limitations.}
\begin{itemize}
\item The intermediate iterates do not represent a common physical time and
cannot describe the timing or physical stability of an oscillation.
\item Internal-face flux contributions need not cancel in the total
inventory. At an exact stationary residual root, the complete cell updates
can vanish, but that does not make the unequal time-weighted contributions
at each face equal. Cancellation can also occur away from a stationary root,
for example when that face's flux is zero.
\item Multiplying each unsplit residual by a positive local interval
preserves its zeros. This argument alone does not establish that the full
finite-interval split chemistry and transport update has the same fixed
points as the stationary residual.
\item Local iterations cannot supply the present acceptance window of three
physical crossing times. They would need to be followed by a global-step
physical march to satisfy the campaign criterion. Stationary certification
is a separate judgment.
\item Convergence or failure of a local relaxation does not establish whether
the physical wind has a stable stationary state or a limit cycle.
\item The implementation passes local intervals to the carrier transport
operator as well as to the hydro update; the performance and fixed-point
effects of that coupling for these campaign states have not been established
by the interval-ratio estimate. \ev{Main}{2309--2311,2608--2609}\ev{D}{2866--2870}
\end{itemize}

Using \texttt{Run mode: init} and local stepping to prepare a He/H ladder
rung before its physical march would be a change to the campaign procedure,
not a measured acceleration or an approved change to the current campaign.

\section*{Sources}

Labels in brackets are file:line references; paths are relative to the
ExoAtmosphere workspace. Values quoted from historical records were not
remeasured by running a solver for this check.
\begin{itemize}
\item M: \path{EXHALE_v1.10/LHS1140b/tools/marching_acceptance.py}.
\item E: \path{EXHALE_v1.10/src/modules/time_step/eval_dt.f90}; IR: \path{EXHALE_v1.10/src/modules/files_IO/input_read.f90}.
\item Par: \path{EXHALE_v1.10/src/modules/init/parameters.f90};
MI: \path{EXHALE_v1.10/LHS1140b/tools/make_march_input.py}.
\item E1: \path{EXHALE_v1.10/LHS1140b/models/k07_conservative_20261030/E1a_time_relaxation_20261104/run/input.inp};
ED: \path{EXHALE_v1.10/md/k07_time_relaxation_design_20261104.md};
CD: \path{EXHALE_v1.10/md/lhs1140b_heh_le14_marching_campaign_20261105.md}.
\item A: \path{EXHALE_v1.10/src/modules/time_step/attempted_step.f90};
G: \path{EXHALE_v1.10/src/modules/init/define_grid.f90};
IO: \path{EXHALE_v1.10/examples/exhale_io.py}.
\item J14: \path{EXHALE_v1.10/LHS1140b/models/marching_campaign_20261105/cases/heh14_s1p60/phase/journal.txt}.
\item J12: \path{EXHALE_v1.10/LHS1140b/models/marching_campaign_three_slots_20261106/cases/heh1p2_s1p28/phase/journal.txt}.
\item J209: \path{EXHALE_v1.10/LHS1140b/models/marching_campaign_three_slots_20261106/cases/heh2p09_s1p42/phase/journal.txt}.
\item J01: \path{EXHALE_v1.10/LHS1140b/models/marching_campaign_he01_series_20261106/cases/heh0p1_s2p53/phase/journal.txt}.
The accepted-step counts are in each phase's compressed chunk log:
\path{chunks/c0049.log.gz}:58299 (J14),
\path{chunks/c0041.log.gz}:63232 (J12),
\path{chunks/c0038.log.gz}:53268 (J209), and
\path{chunks/c0017.log.gz}:68365 (J01); line numbers refer to decompressed text.
\item S12: \path{snapshots/c0041/Hydro_ioniz.txt} in the phase containing J12.
\item Main: \path{EXHALE_v1.10/src/EXHALE_main.f90};
N: \path{EXHALE_v1.10/src/modules/time_step/steady_newton.f90};
D: \path{EXHALE_v1.10/src/modules/lower_atmosphere/diffusive_photochemistry.f90}.
\item K: directory
\path{EXHALE_v1.10/LHS1140b/models/euv_match_20261027/HeH1.2/xuv_x1p28/k07/}.
\item W: directory
\path{EXHALE_v1.00/scratch_codex_q5_q6i_20261106/R1_reference_views/wasp_full_newton/}.
\item P: \path{EXHALE_v1.10/md/lhs1140b_marching_campaign_preflight_20261105.md};
C: \path{EXHALE_v1.10/md/lhs1140b_campaign_representative_s_20261105.md};
Q15: \path{EXHALE_v1.10/md/briefs/q15_he01_packet_20261106_codex.md};
U: project \path{CLAUDE.md}.
\item H: \path{EXHALE_v1.10/LHS1140b/models/k07_conservative_20261030/opcons_20261104/h2p_hehp_same_state_20261104/continuation_20261104/offline_admission_20261104/results/hh_cells_92_105_offline.txt}.
\item Q: \path{EXHALE_v1.10/md/briefs/q13b_carrier_corner_20261106_codex.md}.
\item B: \path{EXHALE_v1.00/md/Update_EXHALE_stage2.md};
L: \path{EXHALE_v1.10/md/briefs/crosscheck_lowheh_composition_20261029_codex_r1.md}.
\item J: \path{EXHALE_v1.10/md/briefs/crosscheck_k07_coupled_gate2c_r3_20261030_codex.md};
JV: \path{EXHALE_v1.10/LHS1140b/models/k07_conservative_20261030/codex_phase_20261101/measured_summary.json}.
\end{itemize}

\end{document}
